\documentclass[11pt]{article}
\usepackage[margin=1in]{geometry}
\usepackage{amsmath,amssymb}
\usepackage{booktabs}
\usepackage{array}
\usepackage{xcolor}
\usepackage{enumitem}
\usepackage{microtype}
\usepackage{titlesec}
\usepackage[hidelinks]{hyperref}
\usepackage{fancyhdr}

\definecolor{good}{RGB}{30,120,60}
\definecolor{bad}{RGB}{170,40,55}
\definecolor{mid}{RGB}{160,120,20}
\newcommand{\g}[1]{\textcolor{good}{#1}}
\newcommand{\bd}[1]{\textcolor{bad}{#1}}
\newcommand{\md}[1]{\textcolor{mid}{#1}}
\newcommand{\code}[1]{\texttt{\small #1}}
\newcommand{\arxiv}[1]{\href{https://arxiv.org/abs/#1}{\texttt{#1}}}

\titleformat{\section}{\large\bfseries}{\thesection}{0.6em}{}
\titleformat{\subsection}{\normalsize\bfseries}{\thesubsection}{0.6em}{}

\pagestyle{fancy}\fancyhf{}
\lhead{\small arxivkg --- phase-5 subgraph retrieval}
\rhead{\small\thepage}
\setlength{\headheight}{14pt}

\title{\vspace{-2.2em}\textbf{Phase-5 Subgraph Retrieval:\\What We Have, What We Tested, What We Don't}}
\author{arxivkg internal analysis}
\date{2026-06-29}

\begin{document}
\maketitle
\thispagestyle{fancy}
\vspace{-1.5em}

\begin{abstract}\noindent
This note inventories the subgraph-retrieval techniques relevant to phase-5 of arxivkg. It
describes the \emph{deployed} method in full (induced-pool + ancestor augmentation + Personalized
PageRank + structural extraction), reports the head-to-head \emph{bake-off} we ran (induce vs.\
PCST/Steiner vs.\ similarity-aware BFS over five topics), and surveys the methods we have
\emph{not} tested, separating those that operate on a \emph{given} graph (in-regime) from those
that \emph{build} a graph from text (a different paradigm). The central empirical finding is that
\emph{membership} and \emph{structure} are separable concerns with different optimal algorithms:
citation-centrality (PPR) maximizes recall of canonical landmarks; cosine-prized Steiner trees
maximize connectivity but discard those landmarks. Our committed recipe takes the best half of
each.
\end{abstract}

\section{The setting, and why it constrains the method space}

Phase-5 turns a query vector into a \textbf{bounded, on-topic, connected citation subgraph} whose
nodes are the ``landmark'' papers of a field and whose edges are the real citations between them.
This subgraph then feeds (a) the phase-6 narrative/compose layer and (b) a lineage-tree
visualization (time on $x$, dendromatic spread on $y$).

The decisive property of our setting is that \textbf{the graph already exists}. We have a real,
directed citation graph: 173M edges in Postgres, each edge carrying an LLM-inferred
\code{intent}$\in$\{extends, uses, compares, critiques, motivates, background\} and a
\code{polarity}. This is categorically different from the mainstream ``GraphRAG'' literature,
where the retrieval graph is \emph{synthesized} from a text corpus by an LLM (entity + relation
extraction) before retrieval can begin. That distinction partitions the entire method space:

\begin{itemize}[leftmargin=1.4em,itemsep=2pt]
  \item \textbf{Group A --- operate on a given graph.} Directly applicable to us; these are
  \emph{selectors} over our citation edges. The bake-off lives here.
  \item \textbf{Group B --- build their own graph from text.} Would have us \emph{discard} the
  citation structure and re-derive a weaker one with an LLM. Not drop-in competitors for the
  retrieval path; at most they donate ideas to an \emph{offline} index/browse layer.
\end{itemize}

A second constraint: phase-6 induces a \emph{small} subgraph ($\sim$400 candidate nodes,
$\sim$5{,}000 edges) and runs everything in Python. We never traverse the 173M-edge graph at query
time, so methods whose entire value proposition is native graph-DB traversal (Neo4j-style) buy us
nothing here.

\section{What we have: the deployed method}

The production path (\code{phases/phase-6/components/\{subgraph,pagerank,recency\}.py}) is a
four-stage pipeline. No LLM and no embedding model run inside it --- the caller passes a precomputed
query vector so VRAM can be sequenced (embed $\rightarrow$ free model $\rightarrow$ load narrator).

\subsection*{Stage 1 --- topical pool (vector + lexical)}
A hybrid candidate pool: pgvector cosine KNN (now on the \emph{asymmetric} nomic space, raw query,
no HyDE --- the phase-4 decision) fused by RRF with a title-BM25 list, then a \emph{frontier pull}
(year-filtered pure-KNN recency top-up, because the corpus has a 2026 delta the topic tree predates).
Result: $\sim$140--510 nodes depending on topic breadth.

\subsection*{Stage 2 --- backward-ancestor augmentation (graph)}
The similarity pool is dominated by recent application/survey papers; the \emph{foundational}
lineage (DDPM, LDM, score-SDE for diffusion) is cited by nearly the whole pool yet is rarely a
top-$N$ embedding neighbour of the query. \code{augment\_with\_ancestors} pulls those most-cited
one-hop ancestors in, ranked by the number of distinct pool papers that cite them, optionally
category-gated so a famous cross-field paper cannot hijack a polysemous query. This is the single
most important recall mechanism --- it is what makes ``Origins'' and the main-path spine exist at all.

\subsection*{Stage 3 --- induce the subgraph (graph)}
\code{induce\_subgraph} keeps every citation edge with \emph{both} endpoints in the pool, collapsing
parallel occurrences to one edge carrying the strongest-signal intent (priority order
extends~$>$~uses~$>$~compares~$>$~critiques~$>$~motivates~$>$~background). Inducing \emph{within} the
similarity-ranked pool --- rather than expanding outward from seeds --- is the key trick that avoids
the hub-explosion trap (ResNet/Adam have $\sim$100k citers): the structure inside the on-topic cloud
is automatically relevant and bounded.

\subsection*{Stage 4 --- rank + structure (Personalized PageRank)}
Raw in-degree favours ancient hubs. Instead \code{rank\_landmarks\_ppr}
(\code{pagerank.py}) runs PPR over the induced edges, \emph{teleporting} back to the
query-similarity weights:
\[
  r \;=\; (1-\alpha)\,p \;+\; \alpha\,\big(W r + (\mathbf{1}^\top_{\text{dangling}} r)\,p\big),
  \qquad \alpha = 0.85,\;\; p_a \propto \max(\text{cos}(a, q),\,0).
\]
This yields \emph{topic-personalized influence}: a paper accumulates rank by being cited
\emph{within this topic's} cloud, so foundational-yet-on-topic papers rise while generic hubs do
not. On top of PPR sit three recency-aware refinements: a topic-adaptive \emph{frontier
reservation} (newest-cohort slots reserved by similarity, since recent papers have low rank by
construction), an \emph{age-normalized} tiebreak (citations-per-year) for the fill, and the
\code{main\_path}/\code{assign\_eras}/\code{tensions} structural extractors that turn the ranked
landmarks into a chronological spine, three eras, and comparison/critique edges.

\medskip\noindent\textbf{In one line:} membership $=$ similarity pool $+$ cited ancestors;
ranking $=$ PPR personalized to the query; structure $=$ longest chronological citation chain.

\section{What we tested: the bake-off}

We held the candidate pool \emph{fixed} (same Stage~1--3 pool, same induced edges, same per-node
cosine ``prize'') and swapped only the Stage-4 \textbf{selector}, across five gold topics
(\code{subgraph\_bakeoff.py}). Three selectors:

\begin{description}[leftmargin=0pt,itemsep=3pt]
  \item[\code{induce} (ours).] PPR-rank the pool to $m$ landmarks. Keeps the whole on-topic pool as
  membership; centrality decides who is a landmark.
  \item[\code{pcst} --- Prize-Collecting Steiner Tree (G-Retriever, \arxiv{2402.07630}).] Find the
  connected subtree maximizing $\sum\text{prize} - \sum\text{edge cost}$. We put the prize on
  \emph{only} the top-$m$ relevant nodes (0 elsewhere) so the solver connects them via minimal
  Steiner \emph{bridges}; edge cost $\approx 0.4\times$ mean prize tunes how many bridges join.
  \item[\code{dsa\_bfs} --- similarity-aware BFS (DynaGRAG, \arxiv{2412.18644}).] BFS outward from
  the top-cosine seeds, gated by $\text{cos}\ge 0.45$, bounded depth~4 / fan-out~8.
\end{description}

\begin{table}[h]\centering\small
\begin{tabular}{l l r r r r r}
\toprule
Topic & Method & $n_{\text{cc}}$ & edges & mean\,cos & recall & note\\
\midrule
diffusion   & \code{induce}  & 20 & 165 & 0.487 & \g{0.625} & best recall, worst connectivity\\
            & \code{pcst}    & 43 & 183 & \g{0.630} & \bd{0.125} & best cos/conn, drops canonical\\
            & \code{dsa\_bfs}& 38 & 344 & 0.587 & \md{0.500} & middle\\
\addlinespace
transformer & \code{induce}  & 28 & 131 & 0.481 & \g{0.600} & \\
            & \code{pcst}    & 40 &  90 & \g{0.592} & 0.400 & \\
            & \code{dsa\_bfs}& 36 & 156 & 0.577 & 0.400 & \\
\addlinespace
testtime    & \code{induce}  & 35 & 204 & 0.522 & \g{1.000} & full canonical recall\\
            & \code{pcst}    & 40 & 189 & \g{0.585} & 0.667 & \\
            & \code{dsa\_bfs}& 33 & 209 & 0.567 & 0.500 & \\
\addlinespace
ssm         & \code{induce}  & 27 & 171 & 0.435 & --- & no gold set\\
            & \code{pcst}    & 26 & 116 & \g{0.556} & --- & \\
            & \code{dsa\_bfs}& 35 & 231 & 0.550 & --- & \\
\addlinespace
gnn         & \code{induce}  & 22 & 146 & 0.554 & --- & no gold set\\
            & \code{pcst}    & 39 & 123 & \g{0.664} & --- & \\
            & \code{dsa\_bfs}& 34 & 227 & 0.646 & --- & \\
\bottomrule
\end{tabular}
\caption{Bake-off ($m\approx40$). $n_{\text{cc}}$: nodes in the largest connected component (higher
$=$ less fragmented). mean\,cos: average query-cosine of chosen nodes (topical-drift proxy). recall:
fraction of the gold canonical set recovered. \g{Green}/\bd{red}/\md{amber} mark the per-column
best/worst/mid where a gold set exists.}
\end{table}

\subsection*{Reading the table}
\begin{itemize}[leftmargin=1.4em,itemsep=2pt]
  \item \textbf{\code{induce} (PPR) wins recall, loses connectivity.} It recovers the canonical
  landmarks (diffusion 0.625, testtime 1.0) because PPR rewards \emph{citation centrality}, and
  canonical papers are central by definition. But the top-$m$ by centrality are not all mutually
  linked, so the induced tree fragments (only 20--35 of 40 in one component).
  \item \textbf{\code{pcst} wins connectivity and mean-cosine, loses recall --- badly.} Because the
  prize is \emph{cosine}, PCST keeps tightly-connected high-cosine clusters and \emph{drops the
  canonical landmarks}, which are often \emph{not} top-cosine neighbours of the query (diffusion
  recall collapses to 0.125). This is the headline finding: \textbf{cosine $\ne$ canonical}. Raw
  KNN recall of canonical papers is only $\sim$0.16; PPR is precisely the fix, and PCST throws it
  away.
  \item \textbf{\code{dsa\_bfs} is a consistent middle} on both axes with no regime where it wins.
  \item An \emph{earlier} PCST run with a uniform low edge cost ($\tau=0.45$) kept the
  \emph{whole} graph (404/468 nodes) --- a degenerate ``connect everything'' solution. The fix
  (prize only on the top-$m$, 0 elsewhere) is what produces the focused $\sim$40-node trees above.
\end{itemize}

\subsection*{The synthesis we committed to}
The two failures are complementary, so the lineage-tree builder
(\code{scratchpad/lineage\_tree.py}) uses \textbf{each algorithm for what it is good at}:
\begin{enumerate}[leftmargin=1.6em,itemsep=2pt]
  \item \textbf{Membership $=$ PPR landmarks} (recall): take the \code{induce} landmark set.
  \item \textbf{Connectivity $=$ PCST with a big prize on \emph{exactly those landmarks}} (0
  elsewhere): the Steiner solver adds the minimal bridge papers that join them into one tree
  ($\sim$36--40 of 40 kept connected).
  \item \textbf{Parent structure $=$ most-recent-older-cited predecessor} (clean layout): each
  paper's parent is the in-set paper it cites with the \emph{latest} year $\le$ its own (tie:
  in-degree). This threads flowing left$\to$right chains (VAE\,$\to$\,U-Net/GANs\,$\to$\,DDPM\,$\to$\,
  Stable~Diffusion\,$\to$\,frontier) instead of the ancient-hub star that raw in-degree parents produce.
\end{enumerate}
Slogan: \textbf{embeddings + PPR decide \emph{who is in}; the graph decides \emph{how it is shaped}.}

\section{What we don't have (yet): surveyed but untested}

The literature survey (run over 5 of the 27 topics) returned the methods below. They split cleanly
along the Group~A / Group~B line from \S1.

\subsection{Group A --- in-regime (operate on our citation graph)}
\begin{description}[leftmargin=0pt,itemsep=4pt]
  \item[SimGRAG (\arxiv{2412.15272}) --- \md{untested, marginal fit}.] Aligns a query to a subgraph
  by graph-\emph{pattern} similarity (a soft subgraph-isomorphism / structural match). Genuinely
  operates on a given graph, so it \emph{is} in-regime. But our queries are \emph{topical}
  (``diffusion models''), not \emph{structural} (``find a motif shaped like $X$''); there is no
  query-side graph pattern to match against. Predicted payoff: low. Worth at most a cheap probe as a
  fourth bake-off selector, mainly to be able to say we tried it.
  \item[Full G-Retriever objective (\arxiv{2402.07630}) --- \md{partially tested}.] We use its PCST
  formulation for connectivity, but \emph{not} its downstream graph-encoder + soft-prompt that feeds
  the subgraph into an LLM via a learned projection. That half is irrelevant to us: our compose layer
  reads the subgraph as text/citations, not as a learned graph embedding. No reason to add it.
  \item[HippoRAG-style PPR (\arxiv{2405.14831}) --- \g{effectively already adopted}.] HippoRAG's core
  idea --- single-step PPR personalized to query-linked nodes over a knowledge graph --- \emph{is}
  what \code{pagerank.py} implements (the module docstring credits it). The piece we don't replicate
  is its synonym/entity-linking front-end (Group~B; not needed when edges are real citations).
\end{description}

\subsection{Group B --- out of regime (build a graph from text)}
\begin{description}[leftmargin=0pt,itemsep=4pt]
  \item[GraphRAG, Microsoft (\arxiv{2404.16130}).] LLM entity/relation extraction $\to$ Leiden
  community detection $\to$ hierarchical community \emph{summaries}, queried for global
  ``sense-making''. Replaces, not uses, the citation graph. \textbf{But} its community-summary idea
  maps onto our retired phase-5 \code{community\_summary} component and is the one Group~B idea worth
  reviving --- as an \emph{offline browse/index artifact}, decoupled from retrieval.
  \item[LightRAG (\arxiv{2410.05779}).] Dual-level (local entity + global theme) retrieval over an
  LLM-extracted KG with incremental updates. Same paradigm mismatch; the ``dual-level'' framing is
  already approximated by our pool (local cosine) $+$ ancestor augmentation (global lineage).
  \item[RAPTOR (\arxiv{2401.18059}).] Recursive clustering $+$ abstractive summarization into a tree;
  purely text-derived, uses \emph{no} edges. Orthogonal to citation-subgraph retrieval. Its recursive
  summary tree could, like GraphRAG, feed an offline topic-browse layer --- not the retrieval path.
\end{description}

\subsection{Methods we are \emph{not} pursuing, and why}
\begin{itemize}[leftmargin=1.4em,itemsep=2pt]
  \item \textbf{GNN re-rankers (GNN-RAG \arxiv{2405.20139} and kin).} Would require training a GNN
  over the citation graph; PPR captures most of the centrality signal at zero training cost and is
  fully interpretable. Not worth the complexity for a retrieval stage.
  \item \textbf{Native graph-DB traversal (Neo4j GDS).} Retired in the store consolidation; our
  subgraph is small and Python-resident, so server-side traversal is moot.
  \item \textbf{Random-walk / node2vec structural embeddings.} We already have semantic (nomic)
  embeddings for membership; a second structural embedding adds a maintenance burden without a clear
  recall or connectivity win the bake-off didn't already get from PPR$+$PCST.
\end{itemize}

\section{Gaps in the evaluation itself}

Honest caveats about how much the bake-off actually proves:
\begin{itemize}[leftmargin=1.4em,itemsep=2pt]
  \item \textbf{Gold sets exist for only 3 of 5 bake-off topics} (diffusion, transformer, testtime);
  ssm and gnn report \code{recall=None}. Recall conclusions rest on three topics.
  \item \textbf{No LLM-judge / end-to-end compose} was run on the subgraph variants --- the bake-off
  is retrieval-level only (recall, connectivity, drift). Whether better connectivity actually yields
  a better \emph{review} is unmeasured. This is the natural next experiment if we want a defensible
  ``why this method'' for the writeup.
  \item \textbf{$m$ is fixed at $\approx$40.} The recall/connectivity trade-off may move at different
  budgets; untested.
  \item \textbf{mean-cosine is a drift \emph{proxy}, not a quality metric} --- a method can score high
  by hugging the query and still miss the field's structure (exactly PCST's failure mode).
\end{itemize}

\section{Recommendation}

\begin{enumerate}[leftmargin=1.6em,itemsep=3pt]
  \item \textbf{Commit the synthesis as the phase-5 default}: PPR membership $+$ PCST-on-landmarks
  connectivity $+$ most-recent-predecessor structure. It is the empirically-justified best-of-both
  and already produces clean lineage trees across all five topics.
  \item \textbf{Keep \code{induce} / \code{pcst} / \code{dsa\_bfs} as config-toggleable selectors}
  (the ablation mechanism), so the phase-5 \code{run.py} can reproduce the bake-off and the writeup
  can cite a real comparison.
  \item \textbf{Treat Group~B as a separate, later, \emph{offline} feature} (community/RAPTOR-style
  summaries for a browse UI), explicitly \emph{not} on the retrieval hot path.
  \item \textbf{Optional one probe before committing}: add SimGRAG as a fourth selector to close the
  ``did you try everything in-regime?'' question; expect it to underperform, then commit.
  \item \textbf{The real missing experiment} is an LLM-judge pass over the subgraph variants
  (does connectivity $\to$ better review?), if we want the phase-5 method choice to be judged on the
  end task rather than on retrieval proxies alone.
\end{enumerate}

\end{document}
